\CourseSection{Testing and Evaluation}\label{sec:testing}

Verification will check metric calculations against independently established reference results and assess data handling, pipeline integration, reproducibility, and performance. Validation will assess whether the measurements support their intended interpretation and the research questions. Each metric owner will document the test inputs, procedure, expected results, acceptance criteria, and associated requirement and prototype IDs.

\CourseSubsection{Metric Verification and Validation}

\begin{singlespace}
\setlength{\LTcapwidth}{\linewidth}
\begin{longtable}{@{}L{0.12\linewidth} L{\dimexpr0.20\linewidth-2\tabcolsep\relax} L{\dimexpr0.68\linewidth-2\tabcolsep\relax}@{}}
\caption{Planned metric verification and validation}\label{tab:testing-plan}\\
\rowcolor{MidGray}
\TableHeader{Test ID} & \TableHeader{Metric} & \TableHeader{Planned verification}\\
\midrule
\endfirsthead
\rowcolor{MidGray}
\TableHeader{Test ID} & \TableHeader{Metric} & \TableHeader{Planned verification}\\
\midrule
\endhead
\bottomrule
\endfoot

T01 & Size
& \textit{Inputs:} 20 skills chosen from the 13,000-skill sample to cover prose-only bodies, bodies with code blocks, bundles with both text and binary files, and at least one skill with \textit{composition\_truncated} = 1. Expected counts come from a short standalone script that tokenizes each tier's text on its own and counts lines and words, and the split between description, body, and bundle is checked by eye. \textit{Procedure:} run pre-processing and the Size module on the 20 skills. \textit{Acceptance:} every count matches the reference script, binary files are left out of the bundle, and every truncated skill is flagged. \textit{Requirements:} FR-01 to FR-04, IF-02, PF-07.\\
\addlinespace[3pt]

T02 & Churn
& \textit{Inputs:} 20 skills chosen from the sample, made up of 10 skills with commit history, 5 edited skills whose repositories are cloned (at least one of them renamed), and one near-clone family of about 5 members. Expected values are worked out by hand from the commit fields and dates, the AI and bot flags are read from the commit messages, the tokens added and removed come from git's word diff of each version and copy pair, and the order of the copies is read from the dates. \textit{Procedure:} run the Churn module on these skills, then have the team label 100 links from the full sample as adapted copy, template only, or unrelated. \textit{Acceptance:} activity values and flags match exactly, token counts are within 2 tokens of git's diff, the history is followed through the rename, each copy in the family is linked to an earlier member, with the direction matching the dates where they differ, and the labelled links meet PF-04. \textit{Requirements:} FR-05 to FR-08, IF-04, IF-05, PF-04, PF-08.\\
\addlinespace[3pt]

T03 & Age
& Compare calculated ages with expected values from known timestamps and a fixed reference date.\\
\addlinespace[3pt]

T04 & Clone coverage
& Compare detected clone groups and reported coverage with a controlled collection whose duplication relationships are known.\\
\addlinespace[3pt]

T05 & Readability
& Compare implemented scores with reference calculations; separately assess whether the scores support the intended readability interpretation.\\

\end{longtable}
\end{singlespace}

\CourseSubsection{Pipeline Verification}

\begin{singlespace}
\setlength{\LTcapwidth}{\linewidth}
\begin{longtable}{@{}L{0.20\linewidth} L{\dimexpr0.38\linewidth-2\tabcolsep\relax} L{\dimexpr0.42\linewidth-2\tabcolsep\relax}@{}}
\caption{Planned pipeline verification}\label{tab:pipeline-verification}\\
\rowcolor{MidGray}
\TableHeader{Test ID} & \TableHeader{Procedure} & \TableHeader{Proposed acceptance criterion}\\
\midrule
\endfirsthead
\rowcolor{MidGray}
\TableHeader{Test ID} & \TableHeader{Procedure} & \TableHeader{Proposed acceptance criterion}\\
\midrule
\endhead
\bottomrule
\endfoot

T06\newline Data handling and integration
& Process controlled records, including missing fields and duplicate content, through preprocessing and the integrated pipeline.
& Every input record has the expected retained, excluded, or flagged outcome; retained identifiers and metric values remain consistent through output generation.\\
\addlinespace[3pt]

T07\newline Reproducibility
& Repeat analysis using the same dataset snapshot, configuration, and software versions. Compare metric tables and values used in plots.
& No differences in deterministic metric values after matching records by identifier; any permitted numerical tolerance is explicitly defined in the corresponding requirement.\\
\addlinespace[3pt]

T08\newline Performance
& Measure elapsed runtime and peak memory for the workload and workstation specified in the performance requirements.
& Runtime and peak memory satisfy the stated acceptance limits for that workload and environment.\\

\end{longtable}
\end{singlespace}

Initial metric verification supports M3, integration testing supports M4, and performance testing supports M5. Output checks support the initial analysis in M6, while interpretation and limitations are reviewed in M10. Applicable tests will be repeated for final validation in M12.

\CourseSubsection{Evaluation and Acceptance}

Test records will identify the dataset or controlled inputs, configuration, expected and actual results, and pass/fail outcome. Failed checks will be investigated and affected tests repeated after corrections. Measurements that fail validation will not support interpretations beyond the evidence available; inconclusive research findings alone do not constitute a test failure.

\clearpage
